\documentclass[10pt,twocolumn,letterpaper]{article}
\usepackage[margin=0.72in,top=0.65in,bottom=0.75in]{geometry}
\usepackage[T1]{fontenc}
\usepackage{times}
\renewcommand{\ttdefault}{lmtt}
\usepackage{microtype}
\usepackage{amsmath}
\usepackage{booktabs}
\usepackage{graphicx}
\usepackage{xcolor}
\usepackage{enumitem}
\usepackage{url}
\usepackage[hidelinks]{hyperref}
\usepackage[font=small,labelfont=bf,skip=3pt]{caption}
\usepackage{titlesec}
\titlespacing*{\section}{0pt}{6pt plus 2pt minus 1pt}{3pt}
\titlespacing*{\subsection}{0pt}{4pt plus 1pt}{2pt}
\titleformat*{\section}{\large\bfseries}
\titleformat*{\subsection}{\normalsize\bfseries}
\setlength{\columnsep}{0.24in}
\setlength{\textfloatsep}{7pt plus 2pt minus 2pt}
\setlength{\dbltextfloatsep}{7pt plus 2pt minus 2pt}
\setlist{nosep,leftmargin=1.1em}
\renewcommand{\arraystretch}{1.05}
\newcommand{\figdir}{../../results/w1/replay/figures/paper}
\newcommand{\hd}[1]{\begin{tabular}[b]{@{}c@{}}#1\end{tabular}}

\title{\vspace{-0.6em}\Large\bfseries Does the Robot Look Done? A Replay Benchmark for\\LLM Supervisors of Learned Robot Policies}
\author{Tom Wang \quad Ranveer Singh \quad Ty Lipscomb\\[2pt]
\small COS 597C, Project 1 report, October 26, 2026\\
\small Code, data, results and exact commands: \url{https://github.com/tom05919/OrcaBench}}
\date{\vspace{-1.6em}}

\begin{document}
\maketitle

\begin{abstract}
\noindent A learned robot policy does not know when it has finished a task, so a language-model supervisor that
watches the cameras has to decide when the task is done and how often to look, and every look costs time. We built a
benchmark that measures these two skills separately and scores them automatically. It replays 168 recorded rollouts
of three learned policies on seven RoboCasa~\cite{robocasa} kitchen tasks. The supervisor's only actions, ``wait $N$ steps'' and
``done'', do not change what the robot does, so a recording is exactly what a live run would show, and every model
sees identical rollouts. A \emph{probe} asks 450 single ``is it done?'' questions labelled by the simulator's success
checker; \emph{supervision sessions} score 168 episodes with a time-charged value against baselines that need no
language model. We evaluated four open vision-language models at three requested reasoning levels (12
configurations). Done-detection balanced accuracy is 0.69--0.78; reasoning raises it by about 0.06 for two models and
not detectably for the other two. No configuration schedules its looks as well as a perfect detector that looks at a
fixed interval, and none has a higher mean value than a fixed timer. A requested reasoning level was not a reliable treatment: for
three of the four models, two levels produced the same amount of reasoning.
\end{abstract}

\section{Motivation}
Vision-language-action policies such as GR00T~\cite{gr00t} and $\pi_{0.5}$~\cite{pi05} carry out household tasks
quickly but unreliably, and they do not stop when they are finished: in our recordings the policy keeps acting until
the step budget runs out, even after the simulator's checker has registered success. A natural remedy is a
supervisor, a language model that periodically looks at the cameras and ends the episode when the task is
done~\cite{saycan,innermonologue}. Whether such a supervisor helps depends on two skills:
\begin{itemize}
\item \textbf{Done-detection:} recognising from camera images that the task is complete.
\item \textbf{Look scheduling:} choosing how long to let the policy run before looking again. Each look costs seconds
on the task clock and money, so looking constantly is not free and looking rarely wastes time after completion.
\end{itemize}
Existing evaluations either test success detection on curated clips~\cite{successdetectors} or report the end-to-end
success of a combined LLM--robot system, where perception, planning and time cost are mixed into one number. Our own
first benchmark did the latter and could not tell which skill failed (Section~\ref{sec:didnotwork}). This report
describes the benchmark we rebuilt to measure the two skills separately, and its results on four models.

\textbf{Contributions.} (1)~A replay benchmark that isolates done-detection (the probe) and look scheduling (the
sessions), scored by the simulator's checker, with baselines that need no language model. (2)~Results for four open
models at three requested reasoning levels, with a check of whether each level changed the reasoning, 95\% cluster
bootstrap intervals, and full accounting of cost and losses. (3)~An account of what did not work.

\section{Benchmark design}
\label{sec:design}
\textbf{Why replay.} Once re-prompting is excluded (it did not help these policies, Section~\ref{sec:didnotwork}),
the supervisor can only say ``wait $N$ steps'' or ``done''. Neither changes the robot, so an uninterrupted recording
of the policy contains everything the supervisor would have seen in a live run. Replay gives every model identical
inputs (no robot randomness between conditions), makes runs repeatable and cheap, and needs no GPU.

\textbf{Data.} 168 recorded rollouts (two machines, 84 each): three policies from the official RoboCasa365 checkpoints (GR00T~N1.5 multitask,
GR00T~N1.5 post-trained on 30\% of the target-task demonstrations, and $\pi_{0.5}$ multitask), seven tasks and four kitchen scenes. Three tasks
are \emph{leave-alone} tasks the policies were trained on (OpenCabinet, PickPlaceCounterToStove, CloseFridge); four
are unseen composite tasks. 57 rollouts reach success. Each rollout has three $256\times256$ camera videos (a frame
every 10 control steps, 20 steps per second) and the checker's value at every step.

\textbf{Probe (done-detection alone).} One question is one model call with no history: the task instruction and
three images (all three cameras now, 2~s earlier and 5~s earlier, tiled). The model answers
\texttt{\{"done": true\}} or \texttt{false}. Positives are taken at the checker's first true frame and 5~s later;
negatives 2~s and 10~s before it, and at the middle and end of rollouts that never succeed (114 positives, 336
negatives, fixed seed). The primary measure is balanced accuracy.

\textbf{Sessions (detection and scheduling together).} A session is one scene's 21 episodes (7 tasks $\times$ 3
policies) in a fixed order, with the policy identified only by an anonymous letter; 8 sessions give 168 episodes per
configuration. In an episode the model starts at step 0 and repeatedly answers \texttt{wait $N$} (10--400 steps) or
\texttt{done}; at each look it sees the cameras now and at two earlier points of the last interval, the step, the time
left, the look cost and its earlier looks. ``Done'' is scored by the checker at that step; reaching the step budget is
a failure. Before each episode it sees one line per earlier episode of the session (policy letter, task, outcome,
looks), and nothing about the current episode's checker.
\emph{A real episode} (Gemma~4, low reasoning, CloseFridge, policy~R): the model waited 100, 200 and 200 steps, then
said ``done'' at step 500; the checker had turned true at step 493, so the episode succeeded with 4 looks and scored
$1 - (25\,\mathrm{s} + 4 \times 2\,\mathrm{s}) / 45\,\mathrm{s} = 0.27$.

\textbf{Score and baselines.} Value $=$ success $\times \max(0, 1 - (\text{robot seconds} + \text{looks} \times c) /
\text{deadline})$, with the task's deadline and $c = 2$~s per look. Three baselines need no language model and pay for
looks in the same way: a \emph{fixed timer} declares at step $T$; a \emph{perfect detector} looks every $k$ steps and
declares at the first look where the checker is true (the bar for scheduling); the \emph{oracle} declares at the
first success. $T$ and $k$ are tuned on the same policy and task \emph{in the other scenes}. The primary
session measure is the model's value minus the perfect detector's on the same recording, on the leave-alone tasks
(72 episodes per configuration, 50 of them on rollouts that succeed).

Table~\ref{tab:design} lists each design decision and the reason for it. All were fixed before the runs they govern.

\begin{table}[t]
\centering\footnotesize
\setlength{\tabcolsep}{3pt}
\begin{tabular}{@{}p{0.33\columnwidth}p{0.625\columnwidth}@{}}
\toprule
\raggedright Decision & \raggedright Reason \tabularnewline
\midrule
\raggedright Replay recorded rollouts, not live runs & \raggedright The supervisor's actions do not change the robot; identical inputs for every model, repeatable, no GPU. \tabularnewline
\raggedright Probe separate from sessions & \raggedright Our live benchmark mixed skills in one score; the probe isolates perception. \tabularnewline
\raggedright Labels from the simulator checker; the 2~s just before success not probed & \raggedright Automatic scoring; the checker lags the images by about 1~s, so frames in that band are ambiguous. \tabularnewline
\raggedright Balanced accuracy & \raggedright Only 25\% of probes are positive, mostly from three tasks; plain accuracy would reward always saying ``not done''. \tabularnewline
\raggedright Same look cost $c = 2$~s for every model & \raggedright Measured latency on a shared API reflects serving load, not the model. \tabularnewline
\raggedright Baselines tuned in other scenes & \raggedright A fixed rule must not be tuned on the rollout it is tested on. \tabularnewline
\raggedright Past outcomes shown, policies anonymous & \raggedright Adapting to a policy is impossible without feedback; the current episode's outcome is never shown. \tabularnewline
\raggedright Same three levels for every model, recorded on every call; 32,768-token reply limit & \raggedright A uniform design; recording the reasoning length checks that a level did something; a high limit avoids penalising long reasoning. \tabularnewline
\raggedright Levels compared within a model only & \raggedright Providers map a requested level onto each model's own settings differently. \tabularnewline
\raggedright 95\% bootstrap clustered by rollout (probe) or task and scene (sessions) & \raggedright Answers about the same rollout are correlated. \tabularnewline
\bottomrule
\end{tabular}
\caption{Design decisions and their reasons.}
\label{tab:design}
\end{table}

\section{Experimental setup}
\textbf{Models.} Kimi K3, Gemma~4 31B, DeepSeek V4.1 Flash and GLM-5.3 Flash: the four models that accept images
through Sail Research's synchronous API, which serves quantized weights. Results are for those deployments.

\textbf{Reasoning levels.} Every model ran at the same three requested levels, \emph{none}, \emph{low} and
\emph{high}, in both the probe and the sessions; GLM rejects disabling reasoning, so its ``none'' is the provider's
\emph{minimal} level. Every call recorded the length of the model's thinking, the stop reason, output tokens and
attempts. Each configuration ran once: 450 probes and 168 session episodes. The grid cost an estimated \$71.9
(Kimi K3 \$58.3); everything run on the API for this study, including runs not analysed here, cost an estimated
\$138.7.

\section{Results}
\label{sec:results}
\begin{table*}[t]
\centering\footnotesize
\setlength{\tabcolsep}{3.5pt}
\begin{tabular}{@{}llrcccrr@{}}
\toprule
& & \multicolumn{3}{c}{Done-detection probe (450 questions)} & \multicolumn{3}{c}{Sessions, leave-alone tasks (72 episodes, 2~s per look)} \\
\cmidrule(lr){3-5}\cmidrule(l){6-8}
Model & Level & \hd{Thinking\\(chars)} & \hd{Balanced accuracy\\{}[95\% CI]} & \hd{Change from none\\{}[95\% CI]} & \hd{Value $-$ perfect\\detector [95\% CI]} & \hd{Looks /\\episode} & \hd{Correct\\(\%)} \\
\midrule
% Generated by scripts/replay_paper_tables.py; do not edit by hand.
Kimi K3 & none & 0 & 0.709 [0.664, 0.755] & --- & $-$0.121 [$-$0.167, $-$0.077] & 10.0 & 76 \\
 & low & 343 & 0.771 [0.733, 0.809] & +0.062 [+0.006, +0.117] & $-$0.064 [$-$0.118, $-$0.012] & 4.6 & 58 \\
 & high & 1,576 & 0.772 [0.727, 0.814] & +0.064 [+0.013, +0.114] & $-$0.071 [$-$0.139, $-$0.023] & 4.4 & 62 \\
\midrule
Gemma 4 31B & none & 0 & 0.757 [0.719, 0.797] & --- & $-$0.061 [$-$0.097, $-$0.028] & 4.3 & 72 \\
 & low & 838 & 0.763 [0.721, 0.804] & +0.006 [$-$0.022, +0.033] & $-$0.085 [$-$0.141, $-$0.032] & 5.5 & 60 \\
 & high & 844 & 0.752 [0.708, 0.793] & $-$0.006 [$-$0.039, +0.026] & $-$0.090 [$-$0.125, $-$0.057] & 5.4 & 60 \\
\midrule
DeepSeek V4.1 Flash & none & 0 & 0.689 [0.633, 0.744] & --- & $-$0.151 [$-$0.221, $-$0.084] & 9.3 & 82 \\
 & low & 1,556 & 0.761 [0.716, 0.805] & +0.073 [+0.021, +0.127] & $-$0.125 [$-$0.182, $-$0.070] & 7.0 & 38 \\
 & high & 1,630 & 0.746 [0.691, 0.797] & +0.058 [+0.006, +0.109] & $-$0.096 [$-$0.155, $-$0.043] & 6.6 & 58 \\
\midrule
GLM-5.3 Flash & none (min.) & 0 & 0.751 [0.705, 0.796] & --- & $-$0.130 [$-$0.197, $-$0.066] & 7.7 & 60 \\
 & low & 0 & 0.753 [0.704, 0.799] & +0.002 [$-$0.043, +0.044] & $-$0.109 [$-$0.166, $-$0.058] & 8.3 & 62 \\
 & high & 206 & 0.784 [0.739, 0.826] & +0.032 [$-$0.007, +0.072] & $-$0.125 [$-$0.183, $-$0.065] & 7.9 & 56 \\
\bottomrule

\end{tabular}
\caption{Results for every configuration. \emph{Thinking}: median characters of the model's reasoning per probe
answer, the check that a requested level changed the reasoning. \emph{Change from none}: paired difference on the same
probes. \emph{Correct}: share of the 50 rollouts that succeed on which the model's ``done'' was
scored a success (the alternative is a premature ``done''). Baselines on the same leave-alone episodes: perfect detector 0.173, timer 0.144, oracle 0.284
(mean value). No reply was cut off by the reply limit.}
\label{tab:results}
\end{table*}

\begin{figure*}[t]
\centering
\includegraphics[width=\textwidth]{\figdir/fig1_probe_accuracy_by_reasoning.pdf}
\caption{Done-detection: balanced accuracy by requested reasoning level, one panel per model (450 probes per point;
95\% bootstrap CI clustered by recorded rollout; unreadable replies count as wrong). The text in each panel is the
paired difference high $-$ none on the same probes. GLM's ``none'' is the provider's minimal level. Levels are
comparable within a panel, not across panels.}
\label{fig:probe}
\end{figure*}

\subsection{Done-detection}
Balanced accuracy lies between 0.69 and 0.78 in every configuration (Table~\ref{tab:results},
Figure~\ref{fig:probe}). Turning reasoning on helps Kimi K3 (high $-$ none $+0.064$ [$+0.013$, $+0.114$]) and
DeepSeek V4.1 Flash ($+0.058$ [$+0.006$, $+0.109$]), mainly because the model says ``done'' more readily when the task is
done: its hit rate rises from 0.64 to 0.85 (Kimi) and from 0.55 to 0.76 (DeepSeek). Gemma~4 does not change although
its reasoning does (0 to about 840 characters), and GLM's only reasoning level adds $+0.032$ [$-0.007$, $+0.072$]. Low
and high never differ. The errors concentrate just before success: every configuration says ``done'' on 62--80\% of
frames 2~s before the robot succeeds, where the task often looks finished, and on 36--50\% of frames from leave-alone
rollouts that never succeed.

\subsection{Look scheduling}
No configuration reaches the perfect detector: all 12 intervals for value minus detector lie below zero
(Table~\ref{tab:results}, Figure~\ref{fig:sessions}). The best configurations (Gemma~4 with no reasoning, Kimi K3 at
low) give up about 0.06 of the detector's 0.173. No configuration has a higher mean value than the fixed timer
(0.144), and only one interval (Kimi K3 high) reaches it.
Reasoning changes the trade-off rather than the outcome. Without it, Kimi K3 and DeepSeek look 9--10 times per episode
and are right when they declare on 76--82\% of the rollouts that succeed, but they pay for every look and declare late
(median 10--18~s after success). With reasoning they look 4--7 times, but declare too early on 38--62\% of the
rollouts that succeed. On the 22 leave-alone rollouts that never succeed, every configuration eventually says
``done''; these episodes score 0 either way. Unseen tasks add nothing: 7 of 96 of those rollouts succeed, and every
value there is at most 0.008.

\textbf{Adaptation.} If a model adapted to the anonymous policies, the policy that usually finishes later should get
the longer first wait and the later declaration. Within task blocks, only one of 24 such comparisons points that way
(Gemma~4 none, first wait, from just 16 untied pairs). For DeepSeek low and all three GLM levels the slower policy was
declared \emph{earlier} more often than not, which fits early false ``done'' on slow policies that fail; we did not
test that explanation. This test was declared exploratory and has
little power (at most 64 policy pairs per configuration).

\begin{figure*}[t]
\centering
\includegraphics[width=0.92\textwidth]{\figdir/fig3_sessions_value_and_tradeoff.pdf}
\caption{Look scheduling on the leave-alone tasks (OpenCabinet, PickPlaceCounterToStove, CloseFridge; 2~s per look).
Left: value minus the perfect detector on the same rollouts, 95\% bootstrap CI clustered by task and scene (0 = as good
as the detector). Right: share of correct declarations on the rollouts that succeed, against looks per episode; the
dashed line joins configurations that no other configuration beats on both axes.}
\label{fig:sessions}
\end{figure*}

\subsection{Is a reasoning level a treatment?}
Only Kimi K3 has three distinct levels (median thinking 0, 343 and 1,576 characters). Gemma~4 and DeepSeek think the
same amount at low and high, and GLM's low level produces no reasoning at all, exactly like its minimal level. A
benchmark that varies ``reasoning effort'' through a provider's setting should therefore measure the reasoning it
actually obtained, as we did, rather than assume the label.

\section{What did not work}
\label{sec:didnotwork}
\textbf{The first benchmark (shown on October 16).} Language-model supervisors (Claude Opus~5.5 and Sonnet~5.5, and
Kimi K3 later) supervised the same three policies live on GPU machines. They could wait, rewrite the policy's
instruction, restart it, or declare done; every model call was charged its measured latency (4.7~s for Opus) on the
task clock, and the headline was regret against a fixed declaration rule tuned in other scenes. Both models lost to
the rule (regret 0.049 [0.004, 0.100] for Opus, 0.066 [0.016, 0.126] for Sonnet) and could not be told apart
($-0.016$ [$-0.058$, $0.028$]). The number did not say \emph{why}, for four reasons:
\begin{itemize}
\item \textbf{One score mixed three skills} (re-prompting, done-detection, look scheduling), so a loss could come from
any of them. The rebuilt benchmark measures them separately.
\item \textbf{Re-prompting could not help.} Even a privileged prompter that reads the simulator's state and gives
correct step-by-step instructions solved 3 of 76 paired unseen-task episodes, against 4 of 76 for the unchanged
policy. Supervisors that rewrote the instruction also failed more often (a small, unpaired comparison). We therefore
removed re-prompting.
\item \textbf{The latency charge decided the result.} In a post-hoc analysis with no charge per call, Opus and Kimi
were statistically tied with the fixed rule. Latency on a shared API measures serving load, so the rebuilt benchmark charges every model the
same 2~s per look.
\item \textbf{Most episodes carried no signal.} The robot alone almost never succeeds on the unseen tasks, so 12 of
21 episodes per session scored about 0 for every supervisor and every rule. Live runs also gave each model different
rollouts, one pass per model, and data spread over three machines.
\end{itemize}

\textbf{Problems during this study, and their fixes.} (1)~The first replay runs sent no reasoning setting, so each
model used the provider's undocumented default (Gemma does not reason by default; the others do), confounding model
with reasoning. Those runs are excluded; the grid above declares and records the level. (2)~A pilot of 10 probes was
used to drop levels that did not change the reasoning; with so few probes it gave models different sets of levels
(and was wrong for DeepSeek), so we replaced it with one grid for every model. (3)~Our spending ledger booked every
timed-out call at its worst case and stopped the runs at \$72 of real spend; timeouts are now counted at the 99th
percentile of real replies, and every call carries an idempotency key so a retry cannot be billed twice. Before the
keys were added, 143 timed-out grid calls may have been billed twice. (4)~Processes crashed 142 times on the
provider's HTTP 404, 400 and 500 errors and on read timeouts, and the computer twice ran out of memory (whole videos
were cached; frames are now decoded on demand). In the last configuration the provider held requests until they timed
out and then reported them lost, and the laptop slept twice with its lid closed; progress fell from about 40 episodes
an hour to 3--4, and we let it finish unchanged. Every look is saved when it arrives and an interrupted episode is
replayed from its start, so every configuration is complete; in the grid, 34 probe replies and 9 session replies were
paid for but lost, 181 episodes were interrupted and replayed (1,402 looks paid twice, \$1.91), and requests in
flight at a crash cannot be counted.

\section{Limitations}
Each configuration ran once, so the intervals cover sampling over rollouts, not run-to-run variation of the model.
Reasoning levels are provider settings, not amounts of reasoning, and the models are quantized deployments. Samples
are small: 57 of 168 rollouts succeed and positives come almost only from three tasks. Supervisors see images only,
every 10 steps. The value score cannot distinguish a false ``done'' from running out of budget on a rollout that never
succeeds. Finally, some ``errors'' just before success may be correct readings of images that the checker labels a
second late.

\textbf{Reproducing.} The repository's README gives the commands; the recorded rollouts and every model reply are
in its release \texttt{replay-data-2026-10-26}, so all tables and figures can be recomputed without API calls.

{\footnotesize
\begin{thebibliography}{9}\setlength{\itemsep}{0pt}
\bibitem{robocasa} S.~Nasiriany et al. RoboCasa: Large-scale simulation of everyday tasks for generalist robots. RSS, 2024.
\bibitem{gr00t} NVIDIA. GR00T N1: An open foundation model for generalist humanoid robots. arXiv:2503.14734, 2025.
\bibitem{pi05} Physical Intelligence. $\pi_{0.5}$: a vision-language-action model with open-world generalization. arXiv:2504.16054, 2025.
\bibitem{saycan} M.~Ahn et al. Do as I can, not as I say: Grounding language in robotic affordances. CoRL, 2022.
\bibitem{innermonologue} W.~Huang et al. Inner monologue: Embodied reasoning through planning with language models. CoRL, 2022.
\bibitem{successdetectors} Y.~Du et al. Vision-language models as success detectors. CoLLAs, 2023.
\end{thebibliography}}

\end{document}
